\chapter{Kết luận và hướng phát triển}

\section{Kết luận}
Dự án đã thay dữ liệu synthetic bằng CrisisMMD thật, sửa lỗi nhãn informative 2.649 mẫu,
kiểm soát duplicate xuyên split và xây dựng quy trình tái lập từ EDA đến quyết định. Sáu
thuật toán được so sánh trên dev sạch. Informative fusion đạt Accuracy 0,7072, F1 0,8079
và F2 0,9035, nhưng F2 chỉ hơn always-positive baseline 0,0096. Humanitarian fusion đạt
Macro-F1 0,4005 so với majority baseline 0,0686 trên test sạch. Giá trị chính không chỉ
là dự báo nhãn, mà là chuyển kết quả thành hàng đợi, mức ưu
tiên, đội tiếp nhận và cơ chế human-in-the-loop có giới hạn năng lực rõ ràng.

Hậu kiểm loại thêm near-duplicate đã xác minh cho F2 0,9006 và Macro-F1 0,3908 mà không
tune lại. Bootstrap cho thấy F2 gain so với dummy có CI 95\% [0,0032; 0,0161], tức khác
0 theo giả định row-bootstrap nhưng vẫn nhỏ về thực tiễn. Biến động theo event và support
rất thấp ở ba lớp khẩn cấp cho thấy hệ thống phù hợp hỗ trợ sàng lọc, chưa phù hợp tự động
ra quyết định sinh tử.

\section{Hướng phát triển}
\begin{itemize}[leftmargin=*]
  \item Tích hợp NER trích xuất địa danh để định vị người kêu cứu.
  \item Crawl dữ liệu thời gian thực; hiệu chỉnh xác suất (calibration).
  \item Thay TF-IDF bằng embedding ngôn ngữ; tinh chỉnh trọng số fusion theo lớp.
  \item Thu thập ground truth Priority và chi phí xử lý để học/tối ưu policy Risk Score.
  \item Leave-one-event-out hoặc group split theo event để đánh giá domain shift nghiêm ngặt hơn.
\end{itemize}
